\begin{solucion}
    Si particionamos el multiconjunto:
    \begin{align*}
        \{1\cdot\infty,2\cdot,\ldots,n\cdot,(n+1)\cdot\infty,\ldots\infty,(n+m)\cdot \infty\}
    \end{align*}
    en dos subconjuntos:
    \begin{align*}
        A=\{1\cdot\infty,2\cdot,\ldots,n\cdot \infty\} \quad \text{y} \quad B=\{(n+1)\cdot\infty,\ldots\cdot\infty,(n+m)\cdot \infty\}
    \end{align*}
    y como para todo $k = 0,1,\ldots,r$, tenemos que $r=k+(r-k)$, entonces para $k \in \{0,1,\ldots,r\}$ calculamos la combinaciones:
    \begin{align*}
       \left| \left(\binom{n+1}{k}\right) \right| &= \left| \{\{k_1 \cdot 1, k_2 \cdot 2, \ldots, k_n \cdot n\}: \sum_{i=1}^n k_i=k\} \right| \\
       &= \left(\binom{n+1}{k}\right)
    \end{align*}
    Asi para el otro subconjunto podemos seleccionar de forma independiente $r-k$ elementos de $B$:
    \begin{align*}
        \left| \left(\binom{m}{r-k}\right) \right| &= \left| \{\{k'_1 \cdot (n+1), k'_2 \cdot (n+2), \ldots, k'_m \cdot (n+m)\}: \sum_{i=1}^m k_i'=r-k\} \right| \\
        &= \left(\binom{m}{r-k}\right)
    \end{align*}
    Por lo tanto, el número total de combinaciones para $r(k)$ es:
    \begin{align*}
       &\left(\binom{n+1}{k}\right) \cdot \left(\binom{m}{r-k}\right) = \\
       &\left| \{\{k_1 \cdot 1, \ldots, k_n \cdot n, k'_1 \cdot (n+1), \ldots, k'_m \cdot (n+m)\}: \sum_{i=1}^n k_i + \sum_{i=1}^m k'_i = r\} \right| \\
    \end{align*}
    Por lo que sumando para $k=0$ hasta $k=m$, pues estos forma casos disjuntos, tenemos:
    \begin{align*}
        \sum_{k=0}^m\left(\binom{n+1}{k}\right) \cdot \left(\binom{m}{r-k}\right) &= \left(\binom{n+m}{r}\right) \\
    \end{align*}
\end{solucion}